\scanpage{10}
III.6

THEOREM III.6. If $\mathcal{L}$ is a locally separable $\mathcal{A}$-algebra, and $\mathcal{R}$ is a commutative $\mathcal{A}$-algebra, then $\mathcal{L}\otimes_{\mathcal{A}}\mathcal{R}$ is a locally separable $\mathcal{R}$-algebra.

\uline{Proof:} We know\figure{\begin{tikzpicture}[baseline=0]
\draw (0,0) .. controls (0.03,-0.18) and (-0.03,-0.43) .. (-0.11,-0.58);
\draw (-0.22,-0.49) -- (-0.11,-0.58) -- (0.01,-0.61);
\end{tikzpicture}} from [4, Corollary 1.6] that $(\mathcal{L}\otimes_{\mathcal{A}}\mathcal{R})_x=\mathcal{L}_x\otimes_{\mathcal{A}_x}\mathcal{R}_x$ is a central separable $\mathcal{R}_x$-algebra for every $x$ in $X$.
\deleted{\illegible{6}}
Thus we need only show that $\mathcal{L}\otimes_{\mathcal{A}}\mathcal{R}$ has finite presentation as an $\mathcal{R}$-algebra. For every $x$ in $X$ there is a neighborhood $U$ of $X$ and an exact sequence
\diagram{\begin{tikzcd}[ampersand replacement=\&]
0 \arrow[r] \& \mathcal{P} \arrow[r] \& \bigoplus^n\mathcal{A}|U \arrow[r] \& \mathcal{L}|U \arrow[r] \& 0
\end{tikzcd}}
where $\mathcal{P}$ is an $\mathcal{A}|U$-module of finite type. Then the sequence
\diagram{\begin{tikzcd}[ampersand replacement=\&]
\mathcal{P}\otimes_{\mathcal{A}|U}\mathcal{R}|U \arrow[r] \& \bigoplus^n\mathcal{A}|U\otimes_{\mathcal{A}|U}\mathcal{R}|U \arrow[r] \& \mathcal{L}|U\otimes_{\mathcal{A}|U}\mathcal{R}|U \arrow[r] \& 0
\end{tikzcd}}
is exact, and $\mathcal{P}\otimes_{\mathcal{A}|U}\mathcal{R}|U$ is an $\mathcal{R}|U$-module of finite type. Thus $\mathcal{L}\otimes_{\mathcal{A}}\mathcal{R}$ is finitely presented as an $\mathcal{R}$-module, and therefore locally separable.

We recall now the notion of the commutant of a subring in a given ring, as described in [4]. This notion, generalized to sheaves, will be of use in the next section.

For any ring $\Lambda$ with center $A$, and any two-sided $\Lambda$-module $M$, $\mathrm{Hom}_{\Lambda^e}(\Lambda,M)$ is naturally isomorphic as an $A$-module with $M^\Lambda$, the $A$-submodule of $M$ consisting of all $m$ in $M$ such that $xm=mx$ for all $x$ in $\Lambda$. If $\Gamma$ is a ring having center $A$, and $\Lambda$ is an $A$-subalgebra

\mspage{9}
